\section{Quality Attributes, Requirements and Testability}
\label{sec:quality-attributes}

\subsection{Modifiability and Extensibility}
\todo{Assess how hard it is to add new cards and expansions (Imploding,
Streaking, Barking Kittens) in the current code. Point to the
\code{if/else} chain over card names and the hard-coded \code{enum Card}.}

\subsection{Testability}
\todo{Summarise testability issues (can reference \cref{sec:unit-testing}).}

\subsection{Other Quality Attributes}
\todo{Pick the relevant ones, e.g. maintainability, readability/cohesion,
coupling, reliability (concurrency in Nope handling), usability of the
text interface, portability.}

\subsection{Quality Attribute Scenarios}
\todo{Write concrete scenarios (source, stimulus, artifact, environment,
response, response measure).}

\begin{table}[H]
\centering
\begin{tabularx}{\textwidth}{@{}lX@{}}
\toprule
\textbf{Part} & \textbf{Modifiability scenario: add a new expansion} \\
\midrule
Source            & \todo{Developer} \\
Stimulus          & \todo{Wants to add the Streaking Kittens expansion} \\
Artifact          & \todo{Card and deck modules} \\
Environment       & \todo{Design time} \\
Response          & \todo{New cards added without modifying existing game logic} \\
Response measure  & \todo{Only new classes added; no changes to existing classes} \\
\bottomrule
\end{tabularx}
\caption{Example quality attribute scenario.}
\label{tab:qas-modifiability}
\end{table}
